% ============================================================
% HEAS — WSC 2026 Submission  (Version 3)
% Track: Modeling Methodology
%
% INSTRUCTIONS FOR COMPILATION:
%   Run: pdflatex paper && bibtex paper && pdflatex paper && pdflatex paper
%   (wscpaperproc.cls and support files are in this directory)
%
% WHAT IS NEW IN V3 vs V2
%   - Abstract: "Three case studies", 1,000-step scale, 32-scenario OOD
%   - §4 Case studies: forward-reference hooks to large-scale results
%   - §5 Evaluation: restructured 4 RQs → 5 RQs
%       RQ2 upgraded: adds large-scale 20-run study (d=1.39, p<10^-6)
%       RQ3 upgraded: adds 32-scenario OOD champion robustness (32/32 wins)
%       RQ4 upgraded: adds scale heatmap + Wolf-Sheep 30-run study
%       RQ5 new:      cross-domain portability (Eco + Enterprise at scale)
%   - 4 new figures: large-scale showdown, noise stability,
%                    champion 32-scenario, scale heatmap
%   - Appendix: 7 new rows for large-scale and OOD results
%   - Conclusion: upgrade summary to match richer evidence base
%
% REQUIREMENT CHECKLIST (WSC 2026):
%   [x] Abstract <= 150 words, single paragraph, no keywords, no math symbols
%   [x] Title bold ALL CAPS (handled by wscpaperproc class)
%   [x] Section titles numbered, bold, ALL CAPS
%   [x] Times New Roman 11pt (handled by wscpaperproc class)
%   [x] 5--12 pages
%   [x] Author biographies mandatory
%   [x] No footnotes
%   [x] No page numbers (handled by wscpaperproc class)
%   [ ] Verify final page count after compile
%   [ ] Fill in real author names/affiliations/emails for camera-ready
% ============================================================

\documentclass{wscpaperproc}
\usepackage{latexsym}      % extra symbols (WSC template standard)
\usepackage{mathptmx}      % Times New Roman for math (WSC template standard)
\usepackage[T1]{fontenc}   % proper font encoding (WSC template standard)
\usepackage{booktabs}      % professional table rules
\usepackage{listings}      % code listings
\usepackage{microtype}     % better text justification
\usepackage{url}           % URL formatting
\usepackage{amsmath}       % math
\usepackage{graphicx}      % figures
\usepackage{pifont}        % \cmark / \xmark
\usepackage{multirow}      % table cell merging
\usepackage[pdftex,colorlinks=true,urlcolor=blue,citecolor=black,anchorcolor=black,linkcolor=black]{hyperref}

\graphicspath{{figs/}}

% ---- Code listing style ---
\lstset{
  basicstyle=\ttfamily\small,
  columns=flexible,
  keepspaces=true,
  frame=single,
  framerule=0.4pt,
  xleftmargin=6pt,
  xrightmargin=6pt,
  breaklines=true,
  showstringspaces=false,
  numbers=none
}

\newcommand{\cmark}{\ding{51}}
\newcommand{\xmark}{\ding{55}}

% ============================================================
\begin{document}

\WSCpagesetup{ZHANG, NIE, and ZHAO}

\title{HEAS: A HIERARCHICAL EVOLUTIONARY AGENT-BASED SIMULATION FRAMEWORK
       FOR MULTI-OBJECTIVE POLICY SEARCH}

\author{\begin{center}Ruiyu ZHANG\textsuperscript{1}, Lin NIE\textsuperscript{2}, and Xin ZHAO\textsuperscript{2}\\
[11pt]
\textsuperscript{1}Dept.~of Politics and Public Administration, The University of Hong Kong, Hong Kong SAR, CHINA\\
\textsuperscript{2}Dept.~of Applied Social Sciences, The Hong Kong Polytechnic University, Hong Kong SAR, CHINA\end{center}
}

\maketitle

\vspace{-12pt}

% ============================================================
\begin{abstract}
HEAS is a reusable framework for \emph{ODE-style and tabular} simulation-based
policy search. It addresses a recurring integration problem in this model
family: the same outcome metric is often recomputed differently across
optimization, policy comparison, and statistical validation. HEAS resolves this
through a uniform \texttt{metrics\_episode()} contract shared by all three
stages, together with a Layer/Stream/Arena composition model, a pluggable
optimizer interface, and a tournament primitive for multi-scenario ranking.
Against a Mesa~3.3.1 baseline, HEAS reduces coupling code from 160 lines to 5,
and from 42 to 5 against a reusable utility-library baseline. Three case
studies---ecological, enterprise, and a Wolf-Sheep mean-field reimplementation
of a published Mesa model---show code reuse across the targeted model family.
In the tested ecological scenarios, the evolved champion wins all 32 held-out
cases and all 24 single-anchor OOD cases.
\end{abstract}

% ============================================================
\section{INTRODUCTION}

Agent-based models are foundational for policy design across ecology,
economics, and organizational science, yet coupling an ODE-style or tabular ABM
to evolutionary optimization requires integration infrastructure that
general-purpose frameworks do not provide. This paper addresses a specific
integration failure: when researchers link an ABM to an evolutionary algorithm,
the same outcome metric is recomputed in incompatible ways across three separate
code paths---optimization, tournament evaluation, and statistical
validation---introducing silent divergence that is difficult to audit. Current
practice resolves this through bespoke coupling code (roughly 160 lines in a
representative Mesa workflow) that entangles metric definition with solver,
comparison, and inference logic.

We introduce HEAS, a framework that standardizes metric definition through a
uniform \texttt{metrics\_episode()} contract consumed identically by the
evolutionary optimizer, tournament scorer, and bootstrap confidence-interval
pipeline. This approach reduces coupling code to five lines and eliminates
conditions for metric divergence.

HEAS targets a specific niche---coupled ODE/tabular dynamics, multi-objective
evolutionary search, and multi-run validation---rather than the full ABM space
including spatially-explicit heterogeneous-agent systems. The framework's four
core components (C1--C4, Section~2) are independently useful and may be adopted
in isolation.

\subsection{The Coupling Code Problem}

Consider a researcher using Mesa to study regulatory policy in a multi-firm
economy. She wants to (1)~evolve a Pareto-optimal tax regime using NSGA-II,
(2)~compare the evolved regime against a reference policy across independent
economic scenarios using tournament voting, and (3)~report bootstrap confidence
intervals. Each step reads the same welfare metric, yet in Mesa each is an
independent code path: the DEAP fitness function extracts welfare from the
model's \texttt{DataCollector} after each run; the tournament scorer extracts
welfare from a separate episode-evaluation function, possibly with different
aggregation; and the bootstrap CI loop extracts welfare from a third data
collection loop, potentially inconsistent with both.

In practice, silent divergence is difficult to detect. Consider an experiment
where the tournament scorer reads the final step's biomass while the
evolutionary algorithm reads the episode mean biomass. The evolved champion may
become the policy with volatile-but-high peaks rather than sustainable mean
production---a consistently applied but undetected inconsistency that could
persist through publication. Such divergences are common precisely because
\emph{no single code path} connects metric definition to all three consumers.

When the researcher adds a second objective---welfare inequality---she must edit
the DataCollector reporter, the DEAP fitness function, the tournament scorer, and
the CI loop, keeping all four in sync manually. A silent divergence can persist
undetected for weeks and invalidate a publication. This pattern, which we call
\emph{coupling code}, appears in representative ABM-based policy-search
workflows~\cite{pangallo2019best,capri2023transport}.

\subsection{Contributions}

This paper makes one primary methodological contribution: it introduces a
validated design pattern that keeps metric definition fixed across optimization,
policy comparison, and statistical inference in ODE/tabular ABM policy search.
The pattern has four separable components:

\textbf{C1} (Uniform contract): a \texttt{metrics\_episode()} scalar-dict
interface consumed identically by the evolutionary optimizer, tournament scorer,
and inference pipeline.

\textbf{C2} (Composition model): a \texttt{Layer}/\texttt{Stream}/\texttt{Arena}
structure that produces C1 without model rewrites.

\textbf{C3} (Pluggable optimizer): a reference NSGA-II wrapper with a replaceable
optimizer interface.

\textbf{C4} (Tournament): a primitive for multi-scenario ranking under four
standardized voting rules.

The paper evaluates this contribution in two ways. First, it shows that the
design sharply reduces coupling code relative to Mesa-based baselines. Second,
it shows that the same metric contract can be reused across optimization,
tournament evaluation, and inference in three case studies without rewriting
framework logic.

% ============================================================
\section{RELATED WORK}

Agent-based modeling is widely used in ecology, economics, and organizational
research~\cite{grimm2006odd,tesfatsion2006ace,bonabeau2002abm}. Frameworks such
as Mesa and NetLogo support model execution, but they do not provide native
infrastructure for evolutionary policy search or statistical comparison across
runs~\cite{macal2010tutorial,kazil2020mesa,wilensky1999netlogo}. The ODD
protocol improves model description and reproducibility~\cite{grimm2006odd}, but
it does not specify how a metric should be shared across optimization,
comparison, and inference.

Simulation optimization has long recognized the gap between solver theory and
simulation-side integration work~\cite{fu2002optimization,hong2009brief}. In ABM
settings, that integration is usually handled through bespoke coupling
code~\cite{pangallo2019best,capri2023transport,zheng2022aieconomist}. Search
libraries such as DEAP, Optuna, and Pagmo provide optimizers, but they do not
define a metric interface reused by optimization, tournament evaluation, and
statistical validation~\cite{fortin2012deap,akiba2019optuna,biscani2020pagmo}.
Ranking-and-selection methods help structure comparison, but are generally
single-objective and do not solve this contract problem~\cite{goldsman1994ranking}.

HEAS is positioned at this interface. It does not propose a new optimizer.
Instead, it standardizes the metric pathway shared by optimization, tournament
scoring, and inference. In that sense it is adjacent to optimizer wrappers,
ensemble managers, and reproducibility tooling, but narrower in scope and more
specific in its methodological target.

% ============================================================
\section{THE HEAS FRAMEWORK}

HEAS is implemented in Python~3.11 and depends on DEAP for evolutionary
algorithms, NumPy for numerical computation, and SciPy for statistical
utilities. All components---EA algorithm, voting rule, statistical
estimator---can be replaced without changing simulation or metric contract code.

\subsection{Layered Composition}

This subsection defines the compositional boundary that lets models be reused
without rewriting evaluation code. The fundamental abstraction is a
\texttt{Layer} with four methods
(\texttt{step}, \texttt{metrics\_step}, \texttt{metrics\_episode},
\texttt{reset}). All downstream components read the same
\texttt{metrics\_episode()} keys.

Layers compose into \textbf{Streams}---named groups forming a coherent
sub-simulation---and Streams compose into an \textbf{Arena}, which orchestrates
step ordering, inter-stream signal routing, and episode-level aggregation. The
naming convention \texttt{stream\_name.metric\_name} ensures globally unambiguous
keys regardless of Arena size. Explicit composition boundaries make components
testable and replaceable without rewriting optimization or tournament code;
the uniform metric contract (C2) remains stable as model internals evolve.

\subsection{Native Multi-Objective EA Integration}

This subsection shows how the metric contract is consumed by optimization
without introducing a separate evaluation interface. The wrapper handles DEAP's creator/toolbox registration, crossover and mutation
wiring, hall-of-fame management, and hypervolume computation. A researcher adds
NSGA-II to a new simulation by writing only the objective function (typically
5--10 lines reading from \texttt{run\_many()} results). The EA algorithm is
replaceable: \texttt{run\_ea()} also accepts CMA-ES and hillclimbing.
Section~\ref{sec:scale} shows that at 1,000-step ecological scale, NSGA-II's
advantage over random search is large ($d\!=\!1.39$), showing that
algorithm-agnosticism and algorithm quality are complementary properties.

\texttt{run\_many()} executes $n$ episodes in parallel via
\texttt{ProcessPoolExecutor}, with per-episode seeds derived deterministically
from a base seed, making individual fitness calls reproducible.

\subsection{Tournament Evaluation}

This subsection defines the comparison layer that uses the same metric contract
for multi-scenario policy ranking. Comparing candidate policies is formalized as a \texttt{Tournament}: a
cross-product of participants and scenarios evaluated for $n_\text{ep}$
episodes and aggregated by argmax, majority, Borda, or Copeland rules. The
rules are intended to reveal whether rankings are robust to aggregation choice.
Universal agreement (as in Section~5.4) indicates clear dominance;
disagreement diagnoses the competitive structure.
HEAS reports rule agreement, sample complexity $P(\text{correct})$ versus
$n_\text{ep}$, and Kendall's~$\tau$~\cite{kendall1938new} stability versus
score noise.

\subsection{Statistical Infrastructure}

This subsection makes the inference layer explicit: the same episode metrics
used by optimization and tournaments also feed uncertainty quantification.
\texttt{heas/utils/stats.py} provides: \texttt{summarize\_runs()} (mean, std,
median, BCa bootstrap 95\% CI~\cite{efron1994bootstrap}),
\texttt{wilcoxon\_test()}, \texttt{cohens\_d()}, and \texttt{kendall\_tau()}
with bootstrap CI. \texttt{heas/utils/pareto.py} provides
\texttt{hypervolume()}~\cite{zitzler1998hypervolume} and
\texttt{auto\_reference\_point()}. All utilities accept \texttt{list[float]}
and return structured dicts, requiring no configuration.

\subsection{Design Constraints and Trade-offs}
\label{sec:design-constraints}

The uniform metric contract \texttt{metrics\_episode()} deliberately implements
a minimalist interface: \texttt{dict[str, float]}. This design choice reflects a
composability-versus-expressiveness trade-off.

\textbf{Minimalism enables composition.} A richer contract---e.g.,
\texttt{dict[str, list[float]]} permitting per-step time series, or
\texttt{dict[str, dict[str, float]]} enabling agent-level distributions---would
require negotiation between model internals and the optimization, tournament, and
inference layers. By constraining the contract to scalar floats, HEAS avoids
entangling aggregation logic with metric definition. A model that produces
agent-level data aggregates once at the \texttt{metrics\_episode()} boundary, then
reuses that scalar summary everywhere downstream. Richer contracts would require
per-stage aggregation choices, reintroducing the coupling-code problem HEAS targets.

\textbf{What is lost.} The scalar contract does not directly represent time-series
outputs or agent-level heterogeneity (e.g., the distribution of firm sizes, not
just mean firm profit). For policy-ranking purposes---the intended use case---this
loss is bounded: selecting among regulatory regimes requires aggregate summaries
as the basis for comparison, even if within-regime agent-level structure remains
unmeasured.

\textbf{NSGA-II as reference implementation.} Section~3.2 adopts NSGA-II as the
default multi-objective evolutionary algorithm based on three criteria:
(1)~established performance on policy-search benchmarks in the simulation
optimization literature; (2)~well-understood parametrization at WSC; and
(3)~ease of replacement via the pluggable \texttt{run\_ea()} interface. The
framework does not depend on this choice: Section~5.4 demonstrates that algorithm
selection should be landscape-aware, and HEAS's agnostic interface enables this
flexibility without code changes.

% ============================================================
\section{CASE STUDIES}

Three case studies span ecological, economic, and published-model domains.

\subsection{Ecological Population Management}

A predator-prey Arena with a 2-gene policy (\texttt{risk},
\texttt{dispersal}) and objectives \texttt{-agg.mean\_biomass},
\texttt{agg.cv}. Five streams model climate forcing, landscape quality, prey
logistic growth, predator dynamics, and cross-stream aggregation; eight
fragmentation$\times$shock scenarios form the evaluation grid.
Section~5.3 validates this model at two simulation scales; the evolved champion
is subjected to a 32-scenario out-of-distribution robustness test in
Section~5.4.

\subsection{Enterprise Regulatory Design}

A four-layer regulatory Arena with a 4-gene policy
(\texttt{tax\_rate}, \texttt{audit\_intensity}, \texttt{subsidy},
\texttt{penalty\_rate}) and objectives \texttt{-welfare.total},
\texttt{welfare.gini}. Firms maximize per-step profit against a stochastic
demand schedule; the regulator, government, and welfare layers implement
enforcement, redistribution, and aggregation. The Pareto trade-off is
aggregate welfare versus distributional equality across 32 governance scenarios.
Section~5.6 reports cross-domain validation at large scale.

\subsection{Wolf-Sheep Predation: Mean-Field Reimplementation}
\label{sec:wolfsheep}

To demonstrate compatibility with \emph{existing published models}, we
re-implement the canonical Mesa Wolf-Sheep Predation model~\cite{wilensky1997wolfsheep,kazil2020mesa}---the flagship example distributed with the Mesa library---as a mean-field ODE Arena.
Its published dynamics are Lotka-Volterra~\cite{lotka1925elements,volterra1926variazioni} predator-prey with grass regrowth on a
20$\times$20 spatial grid.

\textbf{HEAS reimplementation.} We implement the mean-field ODE approximation as a
4-layer Arena, deriving each term directly from Mesa's energy mechanics.
All six published parameters are used verbatim.

\textit{Grass.} Mesa's binary GrassPatch agents collapse to a continuous
density $G\in[0,1]$:
\begin{equation}
  \Delta G = \frac{1-G}{\tau_g} - \frac{S\,G\,\gamma}{N},
  \label{eq:grass}
\end{equation}
where $\tau_g\!=\!30$, $N\!=\!400$, and $\gamma\!\in\![0,1]$ is the policy
grazing-rate gene.

\textit{Sheep} ($g_s\!=\!4$, $r_s\!=\!0.04$):
\begin{equation}
  \Delta S = \Bigl(r_s - \frac{\max(0,\,1-g_s G\gamma)}{g_s}\Bigr)S
             \;-\; \frac{W\,S}{N}.
  \label{eq:sheep}
\end{equation}

\textit{Wolves} ($g_w\!=\!20$, $r_w\!=\!0.05$; $h$ = harvest-rate gene):
\begin{equation}
  \Delta W = r_w \frac{W\,S}{N}
             \;-\; \frac{\max(0,\,1-g_w S/N)}{g_w}\,W
             \;-\; h\,W.
  \label{eq:wolf}
\end{equation}

Every coefficient maps directly to a published Mesa default; no parameters
are ad-hoc. Concretely: \texttt{grass\_regrowth\_time}\,$\to\!\tau_g$,
\texttt{sheep\_gain\_from\_food}\,$\to\!g_s$,
\texttt{sheep\_reproduce}\,$\to\!r_s$,
\texttt{wolf\_gain\_from\_food}\,$\to\!g_w$,
\texttt{wolf\_reproduce}\,$\to\!r_w$,
\texttt{initial\_population}\,$\to\!N$ (density normalizer). The spatial grid
is collapsed to a density field, an acknowledged simplification (Section~6).
Adding EA optimization to this published model required zero additional coupling
code: the 8-line objective function reads \texttt{metrics\_episode()} keys,
identical in structure to the other two case studies.

A 30-run algorithm comparison on this model is reported in Section~5.5,
where the landscape structure determines which optimizer wins.

% ============================================================
\section{EVALUATION}

\subsection{Evaluation Questions and Protocol}

We pre-specify five evaluation questions. Two assess the framework directly:
whether HEAS reduces coupling code (RQ1) and whether the same framework logic
can be reused across domains (RQ5). Three assess whether the framework remains
usable under realistic modeling conditions: ecological search at two scales
(RQ2), tournament and OOD validation in the ecological study (RQ3), and
optimizer comparison on a published Wolf-Sheep reimplementation (RQ4).

These questions are intended to validate framework behavior, not to claim that
HEAS dominates alternative optimizers. The central criterion is stricter: the
same metric definition should flow unchanged through optimization, tournament
evaluation, and inference.

\subsection{Mesa Comparison: Coupling Code Reduction}
\label{sec:mesa-comparison}

We implement the ecological case study as a competent Mesa~3.3.1 model and
count non-blank, non-comment coupling code (model dynamics excluded).
Measured totals: 160~LOC (Mesa) vs.\ 5~LOC (HEAS)---97\% reduction.
To establish a fairer baseline, we additionally implement a reusable
utility-library abstraction (\texttt{mesa\_eco\_util.py}) that wraps
DataCollector setup, metric extraction, episode runners, and seed management
into a single \texttt{MesaEpisodeRunner} class. The call-site coupling code
for this utility-library approach is 42~LOC (measured)---still an 88\%
reduction compared to HEAS's 5~LOC. The utility class itself requires
$\sim$55~additional lines of one-time per-project infrastructure.

\begin{table}[ht]
\caption{Coupling code: Mesa~3.3.1 vs.\ HEAS vs.\ Mesa+utility-library.
  Mesa column: measured from \texttt{mesa\_eco.py}.
  Mesa+Util column: measured from \texttt{mesa\_eco\_util.py} (call-site LOC
  when using a reusable utility class; utility class itself is $\sim$55~LOC,
  a one-time per-project cost). HEAS column: measured from \texttt{heas/eco.py}.}
\label{tab:loc}
\centering\small
\begin{tabular}{p{0.44\linewidth}rrr}
\toprule
Task & Mesa & Mesa+Util & HEAS \\
\midrule
Metric contract / DataCollector setup   & 15 & 8  & 0 \\
Episode metric extraction               & 20 & 6  & 0 \\
Sequential episode runner               & 22 & 4  & 1 \\
Parallel runner (\texttt{ProcessPoolExecutor}) & 20 & 5  & 0 \\
EA fitness glue                         & 14 & 4  & 3 \\
Tournament scorer                       & 18 & 4  & 0 \\
Per-run seed management (30-run)        & 12 & 3  & 0 \\
Bootstrap CI                            & 35 & 6  & 0 \\
Adding a second objective (incremental) & 4  & 2  & 1 \\
\midrule
\textbf{Total}                          & \textbf{160} & \textbf{42} & \textbf{5} \\
\bottomrule
\end{tabular}
\end{table}

Listing~\ref{lst:mesa-blocks} reproduces the two most illustrative
Mesa coupling blocks verbatim, allowing verification without repository
access. Adding a metric requires editing Block~A \emph{and} every downstream
consumer; in HEAS this is a one-line change in one file.

\begin{lstlisting}[caption={Mesa coupling code Blocks A and B (35 LOC from
    \texttt{experiments/mesa\_eco.py}). HEAS equivalent: zero LOC.},
    label=lst:mesa-blocks, language=Python]
# Block A - DataCollector setup (15 LOC)
# Adding a metric: edit here AND the EA function
# AND the tournament scorer, all in sync.
self.datacollector = mesa.DataCollector(
    model_reporters={
        "prey":    lambda m: m.prey,
        "pred":    lambda m: m.pred,
        "extinct": lambda m: float(m.pred <= 0.0),
        "biomass": lambda m: m.prey + m.pred,
    }
)

# Block B - episode metric extraction (20 LOC)
# Every consumer (EA, tournament, CI) calls this
# separately. Silent divergence risk.
def extract_episode_metrics(model):
    df = model.datacollector.get_model_vars_dataframe()
    mean_biomass = float(df["biomass"].mean())
    cv = df["biomass"].std() / mean_biomass \
         if mean_biomass > 0 else 0.0
    return {"mean_biomass": mean_biomass, "cv": cv,
            "extinct": float(df["extinct"].iloc[-1])}
\end{lstlisting}

Table~\ref{tab:loc} shows the full comparison including the Mesa+utility-library
baseline (42~LOC, measured). HEAS's claim over this baseline is that
\texttt{metrics\_episode()} is enforced at the \emph{framework} level, so new
projects inherit the uniform contract without writing or maintaining a
utility module.

\subsection{Multi-Scale Reproducibility}
\label{sec:multiscale}

For RQ2, we test whether the ecological study is reproducible across the
small- and large-scale settings used in this paper. At small scale
($n\!=\!30$, pop\,=\,20, ngen\,=\,10, steps\,=\,150),
ecological HV\,=\,7.665\,(std\,3.518, CI\,[6.424,\,8.914]; HV reference point:
mean\_biomass\,=\,$-$150, CV\,=\,0.5).
The distribution is bimodal, and a $4\!\times\!4$ budget ablation
($n\!=\!20$ runs/cell; S2c) identifies population size as the controlling
factor: escape rate (HV\,$>\!9.0$) rises from 15--25\% at pop\,$\le\!20$ to
40--80\% at pop\,$\ge\!30$.

The study was then replicated at larger scale
(steps\,=\,1,000, n\_eval\,=\,5 episodes per genome evaluation, pop\,=\,30,
ngen\,=\,15, $n\!=\!20$ runs; the reduction from $n\!=\!30$ to $n\!=\!20$
reflects computational cost increase at the 1,000-step horizon, with variance
behavior confirmed robust by the 4\,$\times$\,4 budget ablation study S2c). At this horizon NSGA-II outperforms random
search: mean HV\,=\,16.66\,(std\,10.84,\,CI\,[12.37,\,20.94])
versus random mean\,=\,6.00\,(std\,0.34,\,CI\,[5.85,\,6.14]).
Wilcoxon test: $p\!=\!1.91\!\times\!10^{-6}$; Cohen's $d\!=\!1.39$
(large effect); NSGA-II achieves 178\% higher mean HV than random search.
At 150-step scale the same comparison is not significant ($d\!=\!0.49$,
$p\!=\!0.49$); the Pareto landscape becomes meaningfully richer only at the
longer horizon.

\begin{figure}[ht]
  \centering
  \includegraphics[width=\columnwidth]{fig4_large_scale_showdown}
  \caption{Large-scale algorithm showdown: NSGA-II vs.\ Random at 1,000
    steps/episode ($n\!=\!20$ runs). Bars show mean\,$\pm$\,1\,std; error bars
    show 95\% CI; dots show individual runs.
    Cohen's $d\!=\!1.39$; Wilcoxon $p\!=\!1.91\!\times\!10^{-6}$.}
  \label{fig:large-scale}
\end{figure}

Figure~\ref{fig:large-scale} shows per-run distributions and aggregate
statistics. The HEAS framework required no code changes between the 150-step
and 1,000-step studies: only the \texttt{steps} argument changed.

\subsection{Tournament Validation and OOD Champion Robustness}
\label{sec:tournament}

For RQ3, we evaluate stable tournament behavior and champion performance in the
tested ecological scenario families. Voting-rule agreement and sample complexity
were assessed across 30 repeats of
the 8-scenario tournament with 100 episodes per scenario. All four voting rules
agreed in 100\% of
(scenario, repeat) pairs, and $P(\text{correct winner})\!=\!1.0$ at all tested
episode budgets (4, 10, 25, 50, 100 per scenario). The 100\% agreement likely
reflects clear dominance on these low-dimensional problems: with only 2--4 policy
genes and 8 scenarios, the Pareto landscape offers unambiguous winners. A
near-tie sensitivity study~(S9) confirms that this unanimity disappears when
candidate policies have statistically indistinguishable means.

Noise stability was assessed by injecting synthetic Gaussian noise
$\mathcal{N}(0,\sigma^2)$ into episode scores (30 repeats) and measuring
Kendall's $\tau$~\cite{kendall1938new} between clean and noisy rankings. This
synthetic perturbation serves as a proxy for real stochasticity in outcome
metrics; it characterizes the tournament's mathematical sensitivity to score
noise rather than measuring variance from specific simulation runs. Rankings are stable ($\tau\!=\!1.0$) for noise
up to 0.65\% of the inter-policy margin: $\tau\!=\!0.944$ at $\sigma\!=\!10$
(6.5\% of margin), degrading to $\tau\!=\!0.744$ at $\sigma\!=\!50$.

\begin{figure}[ht]
  \centering
  \includegraphics[width=0.846\columnwidth]{fig8_noise_stability}
  \caption{Kendall's $\tau$ versus injected score noise $\sigma$ ($n\!=\!30$
    repeats, 8 scenarios). Rankings are stable ($\tau\!=\!1.0$) until noise
    exceeds $\sim$6\% of the inter-policy margin.}
  \label{fig:noise}
\end{figure}

OOD champion robustness was validated at two scales. At small scale (16
scenarios, steps\,=\,500), the champion won all 16 ($p\!<\!0.001$). At large
scale (32 scenarios), the champion policy
won all 32 ($\Delta\!=\!+25.0$, $+3.96\%$; $p\!=\!4.66\!\times\!10^{-10}$;
$d\!=\!0.17$). The small effect size ($d\!=\!0.17$) reflects near-deterministic
score distributions in these ecological scenarios; win rate and paired deltas
are the more informative statistics here, as score variance is low across
scenarios. In the single-anchor \emph{true OOD} study (Exp.~C, S6c), the
champion won all 24 scenarios where all four factors lay outside training
($p\!=\!1.19\!\times\!10^{-7}$), matching the 100\% in-distribution win rate.
Here win rate and paired deltas are more informative than Cohen's $d$ because
the score distributions are nearly deterministic.

\begin{figure}[ht]
  \centering
  \includegraphics[width=\columnwidth]{fig7_champion_32scenarios}
  \caption{Champion minus Reference biomass across 32 held-out ecological
    scenarios. Blue bars indicate champion wins (32/32); dashed line marks
    mean $\Delta\!=\!+25.0$ biomass. Wilcoxon $p\!=\!4.66\!\times\!10^{-10}$.}
  \label{fig:champion32}
\end{figure}

\subsection{Algorithm Agnosticism and Scale Sensitivity}
\label{sec:scale}

For RQ4, we test whether framework usefulness depends on any single optimizer.
It does not. In the ecological ablation at steps\,=\,300
(pop\,=\,20, ngen\,=\,10, $n\!=\!10$ runs), Simple hillclimbing outperforms
NSGA-II (Simple\,=\,19.66, NSGA-II\,=\,9.99, Random\,=\,7.61). On this
2-gene landscape, the trade-off is narrow and local search is sufficient. HEAS
therefore remains useful even when the best optimizer is not NSGA-II.

Scale sensitivity is shown in Figure~\ref{fig:heatmap}. Mean HV generally rises
with longer simulation horizons, from 6.0 at 140 steps and 5 episodes per
evaluation to 10.1 at 300 steps and 20 episodes. This indicates that richer
episode horizons produce a more informative search landscape.

\begin{figure}[ht]
  \centering
  \includegraphics[width=0.846\columnwidth]{fig3_scale_heatmap}
  \caption{Scale sensitivity heatmap: mean HV by simulation horizon
    (steps\,$\times$\,episodes/genome). Each cell is $n\!=\!10$ independent
    NSGA-II runs. HV generally increases with steps; non-monotonicities
    between the 300- and 500-step rows reflect high between-run variance.
    Episode count has a weaker secondary effect.}
  \label{fig:heatmap}
\end{figure}

The Wolf-Sheep comparison reaches the same practical conclusion from a different
direction. At steps\,=\,200 (pop\,=\,50, ngen\,=\,25, $n\!=\!30$ runs), Random
achieves the highest HV (14.35), followed by NSGA-II (10.36) and Simple (9.14);
Random also significantly exceeds NSGA-II
($p\!=\!7.99\!\times\!10^{-6}$). The point is not that one optimizer wins in
general. The point is that optimizer performance is landscape-dependent, and
HEAS makes that comparison possible without rewriting framework code.

\subsection{Cross-Domain Portability}
\label{sec:crossdomain}

For RQ5, we ask whether the same framework code can be reused
across domains. The answer is yes at the code level. The ecological and
enterprise studies use the same optimization, tournament, and statistical
infrastructure; only the domain-specific \texttt{metrics\_episode()}
implementations differ. In the enterprise tournament validation across 8
governance scenarios (20 repeats, 30 episodes per scenario), all four voting
rules agree in 100\% of cases and $P(\text{correct winner})\!=\!1.0$ at all
tested budgets. This supports a narrow portability claim: the framework
transfers across the tested ODE/tabular domains without rewriting framework
logic. It does not yet establish portability to spatially-explicit or
heterogeneous-agent ABMs.

\subsection{Threats to Validity}

\textbf{Construct validity.} The case studies are composition vehicles, not
calibrated domain models; results are evidence about framework behavior. The
metric contract prevents \emph{accidental divergence across consumers}: all
three pipeline stages read the same key. An incorrect \texttt{metrics\_episode()}
implementation will be consistently wrong rather than silently inconsistent. The
framework additionally exposes \texttt{validate\_metrics\_episode()}, which raises
\texttt{TypeError} when a \texttt{metrics\_episode()} implementation returns a
non-dict or non-float value, catching contract violations early in development.

\textbf{Internal validity.} Fixed evaluation seeds isolate framework effects
but can suppress variance in some metrics; zero-variance outcomes are treated as
diagnostic signals (e.g., the $K\!=\!120$ degeneracy, enterprise determinism)
rather than performance claims. OOD validation at steps\,=\,500 and on
single-anchor extrapolation scenarios (Exp.~C) extends evaluation beyond the
training horizon, but only within the tested family of ecological scenarios.

\textbf{External validity.} The Mesa baseline is a competent single-project
implementation, not a census of all coding styles. Runtime findings for
\texttt{n\_jobs} depend on episode cost and hardware. All three case studies
use ODE-style or tabular dynamics; researchers working with genuinely
spatially-explicit ABMs---heterogeneous agents on a grid, emergent network
effects, rich move-and-interact loops---will encounter integration patterns
outside HEAS's current design scope. The 97\% LOC reduction is validated
specifically for the tabular policy-search workflow that HEAS targets. The
results therefore support HEAS on the tested model family, not on arbitrary ABM
domains.

% ============================================================
\section{CONCLUSION}

The paper supports a focused claim: HEAS makes simulation-based policy search
more auditable and more reusable by enforcing one metric contract across
optimization, comparison, and inference. The evidence is consistent on this
point. HEAS reduces coupling code sharply relative to Mesa-based baselines,
supports reuse across three case studies in the targeted model family, and
keeps tournament and inference stages on the same metric pathway used in search.

The results also clarify scope. HEAS is not a claim about universal optimizer
superiority. In richer ecological settings, NSGA-II performs strongly. In
simpler landscapes, other search strategies perform as well or better. That is
a feature of the framework rather than a weakness: optimizer choice can vary
without changing framework logic.

\textbf{Limitations.} HEAS is designed for ODE-style and tabular simulations,
not for fully spatial or heterogeneous-agent ABMs. Its metric contract is
scalar-focused (\texttt{dict[str, float]}), so workflows that need time-series or
agent-level distributions require an added aggregation layer. For very light
simulations, process startup overhead can also dominate runtime.

\textbf{Future work.} Priorities are a \texttt{MesaLayerAdapter}, scalable
parallel backends such as Ray or Dask, and automated scenario generation.

For double-blind review, the repository URL is omitted from this version.
Upon acceptance, the full codebase including all experiment scripts and raw
result data will be made available at [repository --- for camera-ready: GitHub
URL + Zenodo DOI]. An anonymized artifact package containing all scripts
under \texttt{experiments/} and raw result files (S1--S10) is prepared for
reviewers.

% ============================================================
\appendix
\section{Reported Means and Sources}

\begin{table}[ht]
\caption{Principal results reported in this paper with source file IDs (S1--S10).}
\label{tab:appendix-results}
\centering\small
\begin{tabular}{p{0.52\linewidth}p{0.23\linewidth}p{0.1\linewidth}}
\toprule
Quantity & Value & Source \\
\midrule
Coupling LOC reduction & 97\% (160 vs.\ 5) & S1 \\
Ecological HV, small scale & 7.665 $\pm$ 3.518 & S2 \\
Enterprise validation HV & 4141.74 / +9.5\% & S2b \\
Noise $\tau$ at $\sigma\!=\!10$ & 0.944 & S3 \\
Noise $\tau$ at $\sigma\!=\!50$ & 0.744 & S3 \\
Simple ablation HV mean & 19.66 & S4 \\
NSGA-II ablation HV mean & 9.99 & S4 \\
Random ablation HV mean & 7.61 & S4 \\
\midrule
Ecological HV, large scale & 16.66 $\pm$ 10.84 & S5 \\
Random HV, large scale & 6.00 $\pm$ 0.34 & S5 \\
NSGA-II vs.\ Random $d$ & 1.39 & S5 \\
16-scenario OOD wins & 16/16 ($p\!<\!0.001$) & S6a \\
32-scenario OOD wins & 32/32 / +25.0 (+3.96\%) & S6b \\
\textit{True OOD} 24-scenario wins & 24/24 / $1.19\!\times\!10^{-7}$ & S6c \\
Wolf-Sheep NSGA-II vs.\ Random $p$-value & $7.99\!\times\!10^{-6}$ & S7 \\
Enterprise portability HV & $5.565\!\times\!10^{6}$ & S8 \\
Near-tie mean off-diagonal agreement & 0.755 (24.5\% disagree) & S9 \\
Enterprise tournament agreement (100\%) & 100\% / $P\!=\!1.0$ & S10 \\
\bottomrule
\end{tabular}
\end{table}

\noindent\textbf{Source IDs} (paths under \texttt{experiments/results/}):

\begin{description}\setlength{\itemsep}{0pt}\setlength{\parsep}{0pt}\small
  \item[S1] \nolinkurl{mesa_vs_heas/loc_comparison.json}
  \item[S2] \nolinkurl{eco_stats/pop20_ngen10_trait/summary.json}
  \item[S2b] \nolinkurl{ent_stats_small_n20/summary.json}
  \item[S2c] \nolinkurl{bimodal_ablation/ablation_summary.json}
  \item[S3] \nolinkurl{tournament_stress/noise_stability.json}
  \item[S4] \nolinkurl{baseline_comparison/ablation.json}
  \item[S5] \nolinkurl{large_scale_comparison/p1_summary.json}
  \item[S6a] \nolinkurl{baseline_comparison/champion_vs_ref.json}
  \item[S6b] \nolinkurl{large_scale_comparison/p3_summary.json}
  \item[S6c] \nolinkurl{true_ood/phase2_summary.json}
  \item[S7] \nolinkurl{wolf_sheep_study/algo_comparison.json}
  \item[S8] \nolinkurl{large_scale_comparison/p2_summary.json}
  \item[S9] \nolinkurl{near_tie_tournament/best_case.json}
  \item[S10] \nolinkurl{enterprise_tournament/summary.json}
\end{description}

% ============================================================
\begingroup
\small
\bibliographystyle{wsc}
\bibliography{references}
\endgroup

% ============================================================
\section*{AUTHOR BIOGRAPHIES}

\noindent {\bf \MakeUppercase{Ruiyu Zhang}} is affiliated with the Department of Politics and Public Administration at The University of Hong Kong, Hong Kong SAR, China. Her e-mail address is \email{Ruiyuzh@connect.hku.hk}.\\

\noindent {\bf \MakeUppercase{Lin Nie}} is affiliated with the Department of Applied Social Sciences at The Hong Kong Polytechnic University, Hong Kong SAR, China. His e-mail address is \email{lin-apss.nie@polyu.edu.hk}.\\

\noindent {\bf \MakeUppercase{Xin Zhao}} is affiliated with the Department of Applied Social Sciences at The Hong Kong Polytechnic University, Hong Kong SAR, China. His e-mail address is \email{xinnn.zhao@connect.polyu.hk}.

\end{document}
